\chapter{Introduction}

Computational modelling of physical systems plays a central role in many scientific disciplines and engineering applications. By simulating real-world systems through computational models, researchers can perform experiments that would be impractical or prohibitively expensive at full scale. It also allows for making more informed predictions on how the modelled system would evolve in the future, informing decision-making across various domains. Due, arguably, to its profound influence on human life, weather has long been a focus of mathematical modelling, with early numerical prediction approaches dating back to the early 20th century\cite{primitiveeqs, richardson1922weather}. As computational resources have expanded, atmospheric modelling has advanced in parallel. Traditional approaches, most notably the numerical solutions of governing equations, have evolved into the main backbone of modern weather forecasting centers. However, faithfully capturing small-scale, often chaotic dynamics requires physics-based models to run high-fidelity simulations with fine spatial and temporal resolution, resulting in high computational costs and slow predictions \cite{ECMWF2025_HPC_weather}.

Recent progress in machine learning (ML) algorithms and hardware acceleration has enabled a novel line of physical modelling, based purely on data \cite{rosenblatt1957perceptron, raissi2017physics, scher2018toward}. These methods learn the dynamics of the modelled system directly from observations and typically require less explicit domain expertise than equation-based models. In several problem setups, such as global atmospheric modelling, data-driven approaches have achieved performance that rivals or even surpasses state-of-the-art numerical systems while demanding substantially fewer computational resources \cite{aurora, graphcast, fourcastnet}. Despite its success, there are important limitations to this approach. Firstly, training high-capacity ML models often requires very large, high-quality datasets that may be unavailable or costly to curate. Secondly, the black-box nature of such models makes it difficult to uncover their inference mechanisms. While this may not matter when the goal is simply producing accurate predictions, it limits the model’s ability to uncover interpretable dynamics. As a result, such models are less useful for tasks that require hypothesis generation, hypothesis testing, or a deeper understanding of the underlying weather processes. Thirdly, because such models do not explicitly encode physical structure, they are often harder to generalise to unseen settings, such as longer time horizons or different climate regimes.

To address the shortcomings of purely numerical and purely data-driven approaches, hybrid modelling frameworks have been proposed \cite{chen2018neural, neuralgcm}. The core idea is to couple a numerical solver with a learned component, so that the physical structure and the domain priors are incorporated into the model setup, while latent, often sub-scale information is captured from empirical observations. A common practice for implementation is to equip the solver with a learned module, which during the simulation corrects the numerical predictions to minimize some error metric derived from the numerical predictions and the empirical observations.

Hybrid methods offer several benefits: Compared to purely neural solvers which are physics agnostic, they allow for more accurate physical representation and provide trajectories in physical space. They also exhibit improved stability over long timeframe predictions, and require singificantly less empirical data to train on. Compared to purely numerical approaches, they offer significant speedups \cite{Kashinath2021, Meng2025, Kochkov2021}. Thus, it is not surprising, that several such models have been applied to physical simulations, such as fluid dynamics \cite{Kochkov2021}, heat propagation \cite{Liao2023} and atmospheric modelling \cite{neuralgcm}, where they achieve strong performance.

There is another significant advantage hybrid models exhibit, to which I would like to devote the following paragraphs, which is their ability to adhere to closure problems. These problems are apparent in many physical and mathematical domains, and arise because certain variables of interest depend on quantities that are not resolved within the model. For example, while the main physical laws, from which the dynamic evolution of the entire system can be derived are known, certain processes take place on spatial and temporal scales, that in practice are simply impossible to represent within the model, even though they significantly affect the large scale quantities. If these small-scale dynamics are not resolved, the large-scale systems are unclosed and we are facing a closure problem \cite{sanderse2024scientific}. 

For example, in global weather forecasting, deep moist convection often occurs on scales that are too small to be fully resolved by the computational grid. Yet these unresolved processes strongly affect resolved large-scale quantities through, for instance, latent heating, momentum transport, and the organization of precipitation systems. As a result, their aggregate effect must be represented through subgrid parameterizations, making atmospheric convection a classical closure problem \cite{wedi2020baseline}. In large-eddy simulation of turbulent combustion, the resolved equations capture the large turbulent structures, but smaller turbulent motions, scalar transport, and parts of the turbulence–chemistry interaction remain unresolved. \citet{janicka2005large} note that the filtered reacting-flow equations contain unknown subgrid terms such as the SGS stress tensor, SGS species and heat transport, and the filtered chemical source term, all of which require closure models. In cosmological galaxy-formation simulations, black-hole accretion and feedback occur on scales far below the grid resolution. \citet{negri2017black} emphasize that the relevant accretion scale, the Bondi radius, is usually unresolved, so cosmological hydrodynamic simulations resort to sub-grid prescriptions for black-hole growth and AGN feedback.

Constructing an adequate closure model is notoriously hard and traditionally had been done by approximating the effects of small-scale processes on large-scale quantities via some form of forcing. Since these mappings are strongly non-linear, they traditionally have been based on domain specific heuritics, and even then were mostly a game of trial and error, rather than science. With the advent of neural networks, which are known to be strong approximators of non-linear mappings \cite{hornik1989multilayer}, there is growing momentum of parametrising the closure term by neural networks and learning the mappings by optimising the network weights using data. This approach is termed neural closure modelling and has been successfully applied to many problems, as discussed in Section \ref{LR} \cite{Gupta2021}.

However, while applications of neural closure models are successful, the dynamical properties of such systems have not been extensively studied. In particular, the interpretability of learned closures is still a major open issue, especially when the closure is represented by a black-box network \citep{Sanderse2025,Beck2021}. As a result, it is often unclear which processes a learned closure is actually capturing. For instance, the closure forcing may account for missing physical effects, leading numerical-error terms, or unresolved subgrid transfer processes. Relatedly, it often remains unclear on which spatiotemporal scales these corrections interact with the resolved dynamics \citep{Gupta2023,Portwood2021}. This gap in understanding is partially due to the inherent black-box nature of neural networks, but also the complexity of analysing coupled learned–mechanistic dynamical systems. In realistic multiscale settings, it is rarely possible to derive theoretical characterisations of learned closures that remain both mathematically precise and practically relevant. As a result, most existing theoretical analyses have been restricted to simplified textbook cases which, while insightful, have limited applicability to real-world problems. Moreover, even from an empirical perspective, methods for interrogating learned closures in realistic systems remain limited. Much of the literature still emphasizes a priori and a posteriori predictive performance, while the diagnostics of learned closures remain comparatively less common. Existing interpretability methods have yielded useful insights, but they have mostly been demonstrated in idealized turbulence or simplified ocean settings, with fewer comparable analyses in realistic forecasting systems.

To address this gap, we pursue an empirical approach based on a hybrid model in a realistic global weather forecasting setting. Since real-world multiscale systems are highly complex and may involve unresolved or only partially understood processes, deriving general theoretical characterisations of learned closures is often infeasible in practice. We therefore focus instead on empirical analysis of how the learned closure interacts with the resolved dynamics. Our aim is to use diagnostic tools that are familiar within the dynamical-systems, numerical-analysis, and harmonic-analysis communities, and that do not rely on interpreting the internal representation of the neural network itself. A key requirement for this programme is that the full hybrid model be differentiable, since the diagnostics we employ depend on derivatives of the learned component, such as Jacobians. To this end, we implement the model in the \texttt{jax} framework, which enables efficient automatic differentiation through the full simulation pipeline \cite{wengert1964simple,linnainmaa1976taylor,jax}.

Concretely, we develop a fully differentiable hybrid forecasting model that delivers competitive predictions i the global weather forecasting setting also used by operational models, while also enabling empirical analysis of the learned closure. The model is constructed using the neural differential equation formulation and combines a spectral dynamical core with a GNN-based closure \cite{neuraldiffeqs2}. It is trained end-to-end in an online autoregressive manner to predict seven atmospheric variables. During simulation, we quantify the scale selectivity of the learned component by projecting its Jacobian onto an orthogonal basis, which then allows us to study how the closure interacts with the resolved dynamics across forecast horizons and grid resolutions.

While the main architectural choices are discussed in detail in Chapter \ref{Methods}, we briefly motivate them here because they are central to the analyses carried out later. First, we adopt a spectral dynamical core, since spectral methods are widely used in atmospheric modelling on the sphere. In particular, spherical harmonics provide a natural orthogonal basis for representing fields on the globe and enable efficient spectral transforms with high accuracy. Second, we use a GNN-based closure for two reasons. On the one hand, GNNs are well suited to modelling local interactions on irregular discretisations. On the other hand, our scale analysis also requires an orthogonal basis, for which graph Laplacian eigenmodes provide a natural analogue. This makes the proposed cross-scale analysis adaptable beyond the present application and, in principle, transferable to other settings in which a GNN is used as a closure model.

This thesis is organized into 6 chapters. Chapter~\ref{LR} establishes the theoretical framework for neural closure modeling and Graph Neural Networks (GNNs) in atmospheric modeling. Chapter~\ref{Methods} presents our neural closure model, detailing the spectral numerical core, the GNN-based neural closure term and the ground-truth dataset. It also formalizes the scale–interaction experiments conducted with the model. In Chapter~\ref{AER}, we evaluate the model’s predictive performance and compare the performance of two training methodologies (offline training vs.\ trajectory fitting). We present and analyze the interaction types and scales exhibited by the hybrid components and compare them with the ground truth. Finally, Chapter~\ref{Discussion} interprets the results, and Chapter~\ref{CAFW} discusses implications for future research and concludes the thesis.
